\chapter{Chapter06 — Deterministic Query}
%%% generated by the blueprint pipeline from FABL.Chapter06.LearningAndTesting.DeterministicQuery
\section{Overview}

Declarations of the Lean module \texttt{FABL.Chapter06.LearningAndTesting.DeterministicQuery}.

\section{Declarations}

\begin{definition}[Deterministic query programs]
\lean{FABL.DeterministicQueryProgram}
\label{FABL.DeterministicQueryProgram}
\leanok
Given query, answer, and output types, a deterministic query program is built from three
forms: returning a specified output; querying a finite batch of inputs and continuing
with a program determined by the answers to the entire batch; or charging a specified
number of local computation steps before continuing.
\end{definition}

\begin{definition}[Query batch cost]
\lean{FABL.DeterministicQueryProgram.queryBatchCost}
\label{FABL.DeterministicQueryProgram.queryBatchCost}
\leanok
\uses{FABL.LearningCost}
For a batch of $m$ queries, define its cost to have zero random examples, $m$ membership
queries, and $m$ charged local computation steps.
\end{definition}

\begin{definition}[Program bind]
\lean{FABL.DeterministicQueryProgram.bind}
\label{FABL.DeterministicQueryProgram.bind}
\leanok
\uses{FABL.DeterministicQueryProgram}
Given a rule assigning a deterministic query program to each output, bind a
deterministic query program to that rule by replacing each pure output with the program
assigned to it, preserving each query batch and its answer-dependent continuation, and
preserving each charged computation step and its continuation.
\end{definition}

\begin{definition}[Output map for query programs]
\lean{FABL.DeterministicQueryProgram.map}
\label{FABL.DeterministicQueryProgram.map}
\leanok
\uses{FABL.DeterministicQueryProgram, FABL.DeterministicQueryProgram.bind}
Given types of queries, answers, outputs, and results, a function from outputs to
results induces a query program with results as its output type: it runs the original
program and applies the function to the output it returns.
\end{definition}

\begin{definition}[Execution with resource cost]
\lean{FABL.DeterministicQueryProgram.runWithCost}
\label{FABL.DeterministicQueryProgram.runWithCost}
\leanok
\uses{FABL.DeterministicQueryProgram, FABL.DeterministicQueryProgram.queryBatchCost, FABL.LearningCost}
Given types of queries, answers, and outputs, and an oracle assigning an answer to each
query, every deterministic query program evaluates to an output together with a cost
triple. A pure output incurs zero cost; a batch of $m$ queries adds $m$ membership
queries and $m$ units of local work; a charged computation step adds its specified
number of units of local work. The random-example cost is zero throughout.
\end{definition}

\begin{theorem}[Cost of sequential query programs]
\lean{FABL.DeterministicQueryProgram.runWithCost_bind}
\label{FABL.DeterministicQueryProgram.runWithCost_bind}
\leanok
\uses{FABL.DeterministicQueryProgram, FABL.DeterministicQueryProgram.bind, FABL.DeterministicQueryProgram.runWithCost}
\difficulty{4}
Let a deterministic query program return an output $x$ with cost $c$, and let a
continuation assign a query program to each possible output. Running the program
obtained by sequencing this continuation after the first program returns the
continuation's output and has cost equal to $c$ plus the cost of running the
continuation at $x$.
\end{theorem}

\begin{definition}[Conversion to a query-learning program]
\lean{FABL.DeterministicQueryProgram.toLearningProgram}
\label{FABL.DeterministicQueryProgram.toLearningProgram}
\leanok
\uses{FABL.DeterministicQueryProgram, FABL.LearningProgram, FABL.LearningProgram.bind, FABL.Sign, FABL.queryInputBatchProgram}
For each dimension $n$ and output type, every deterministic program whose queries are
points of the $n$-dimensional Boolean cube and whose answers are signs can be
represented as a query-learning program with the same output type. Each finite batch of
inputs is queried against the target, and the program continues using the resulting
labels; charged local computation steps are retained.
\end{definition}

\begin{theorem}[Deterministic programs as learning programs]
\lean{FABL.DeterministicQueryProgram.runWithCost_toLearningProgram}
\label{FABL.DeterministicQueryProgram.runWithCost_toLearningProgram}
\leanok
\uses{FABL.BooleanFunction, FABL.DeterministicQueryProgram, FABL.DeterministicQueryProgram.runWithCost, FABL.DeterministicQueryProgram.toLearningProgram, FABL.LearningProgram.addOutcomeCost, FABL.LearningProgram.runWithCost, FABL.LearningProgram.runWithCost_bind, FABL.Sign, FABL.runWithCost_queryInputBatchProgram}
\difficulty{4}
Let $n$ be a natural number, let the target be a function from the $n$-dimensional sign
cube to the sign values, and let a deterministic query program query points of that cube
and return an output of any type. Interpreting it as a learning program that makes
membership queries yields a point-mass law concentrated on the deterministic run's
output and cost.
\end{theorem}
